% ============================================================================
%  CHƯƠNG 1. GIỚI THIỆU
% ============================================================================
\section{Giới thiệu}
\label{sec:intro}

\subsection{Bối cảnh}

Ung thư phổi là nguyên nhân tử vong do ung thư hàng đầu thế giới~\cite{bray2024globocan}. X-quang ngực
(CXR) là xét nghiệm hình ảnh được chỉ định nhiều nhất, và nhiều ca ung thư phổi lần đầu lộ ra dưới dạng
\emph{nốt} (đường kính $\le 3$\,cm) hoặc \emph{khối} ($> 3$\,cm) trên CXR. Nốt nhỏ, mờ hoặc bị xương sườn,
tim che khuất rất dễ bị bỏ sót: trên các tập làm giàu ca bệnh, sensitivity của người đọc chỉ khoảng
57--65\%~\cite{sim2020dcnn,ahn2022aiaided}, và với các ca ung thư từng bị bỏ sót trước đó chỉ còn
40--53\%~\cite{jang2020dlad}.

Các hệ thống CAD dựa trên CNN đã đạt AUROC 0,92--0,99 trên các tập ngoài~\cite{nam2019dlad}, và một thử
nghiệm ngẫu nhiên có đối chứng cho thấy CAD làm tăng tỷ lệ phát hiện nốt cần xử trí từ 0,25\% lên
0,59\%~\cite{nam2023rct}. Song song, các mô hình ngôn ngữ--thị giác (VLM) và LLM đa phương thức đang được
giới thiệu như \emph{trợ lý đọc CXR}: chúng nhận ảnh, trả lời bằng ngôn ngữ tự nhiên và có thể giải thích.
Câu hỏi thực tế đặt ra là: các trợ lý này có thật sự \emph{thấy} nốt phổi hay không?

\subsection{Vấn đề}

Các đánh giá độc lập năm 2026 cho câu trả lời đáng lo ngại:
\begin{itemize}
  \item Sáu VLM zero-shot trên 247 ảnh JSRT chỉ đạt sensitivity 0,6--27,3\%. MedGemma-4b-it đạt 5,2\%
        (8/154) ở specificity 100\%; với nốt mờ nhất (độ mờ 1), mọi mô hình đều $\le 12\%$~\cite{nishio2026zeroshot}.
  \item Chín hệ thống trên MosMedData VII (83 ảnh): AUC tốt nhất 0,743 (BiomedCLIP); LLM đa năng tốt nhất
        (Claude 3.7 Sonnet) đạt AUC 0,625; CXR-LLaVA chỉ đạt sensitivity 0,079~\cite{khovanova2026comparative}.
        Trên cùng nguồn dữ liệu, CAD CNN thương mại đạt AUC 0,956~\cite{arzamasov2024independent}.
  \item Bốn MLLM đa năng chỉ phát hiện khoảng 10\% ca ác tính đơn ổ~\cite{chen2026limitations}; thêm ngữ
        cảnh lâm sàng có cấu trúc giúp accuracy chung tăng, nhưng ung thư phổi vẫn chỉ đúng 20\%~\cite{cusick2026context}.
\end{itemize}

Kiểu sai phổ biến ở các VLM chuyên CXR là \textbf{chế độ báo bình thường} (\emph{normal-default mode}):
specificity gần 1, sensitivity gần 0. Có hai giả thuyết bổ trợ nhau. Về thị giác, bộ mã hoá ViT làm mất tín
hiệu tổn thương nhỏ ở bước gộp toàn cục; probe tuyến tính trên patch cục bộ lấy lại được tín hiệu, nâng AUC
nốt của MedSigLIP từ 0,805 lên 0,955~\cite{muthyala2026frozen}. Về dữ liệu, báo cáo huấn luyện phần lớn là
bình thường, nên mô hình nghiêng về câu trả lời ``không có bất thường'' --- một dạng \emph{học lối tắt} khi
lớp dương hiếm.

Đã có tiền lệ tinh chỉnh VLM/LLM trên CXR, nhưng cho tác vụ khác: QLoRA một VLM nhỏ để sinh báo
cáo~\cite{volkov2025lora}, LoRA cho sinh báo cáo~\cite{gao2026s2d}, và tinh chỉnh Qwen2.5-VL cho grounding
cụm từ tổng quát~\cite{xu2025medground}. Trong phạm vi 90 bài nhóm đã khảo sát, \textbf{chưa có nghiên cứu
nào tinh chỉnh VLM cỡ nhỏ để phát hiện và định vị nốt/khối phổi trên CXR rồi so trực tiếp với CNN trên cùng
cách chia dữ liệu}, và cũng chưa có nghiên cứu nào kiểm tra mô hình sau tinh chỉnh có nhìn đúng vị trí nốt hay không.

\subsection{Mục tiêu}

Đồ án xây dựng và đánh giá một \textbf{nguyên mẫu trợ lý AI} nhận một ảnh CXR, cho biết có nốt/khối phổi hay
không và chỉ ra vị trí bằng bounding box. Trợ lý dựa trên VLM cỡ nhỏ (3--4 tỷ tham số) tinh chỉnh bằng LoRA,
đủ nhẹ để huấn luyện trên một GPU Colab. Cụ thể:
\begin{enumerate}
  \item \textbf{Kiểm toán zero-shot có chuẩn hoá:} đánh giá MedGemma-4B và Qwen2.5-VL-3B với prompt chung,
        đầu ra có cấu trúc, ngưỡng chọn trên validation và chạy lặp, đặt cạnh baseline CNN trên cùng split.
  \item \textbf{Tinh chỉnh LoRA} cho hai bài toán: (a) có hay không có nốt/khối; (b) có nốt/khối kèm box.
  \item \textbf{Chống chế độ báo bình thường:} so sánh LoRA chuẩn với LoRA cân bằng lớp.
  \item \textbf{Kiểm tra trợ lý có nhìn đúng chỗ không:} định vị (IoU, hit rate), phép thử che nốt so với che
        vùng ngẫu nhiên cùng diện tích, phân tầng theo kích thước và độ mờ của nốt.
\end{enumerate}

\begin{keybox}
\textbf{Câu hỏi nghiên cứu chính.} (Q1) LoRA có đưa sensitivity của VLM, ở specificity 0,90, lên tối thiểu
0,60 trên tập test trong phân phối không? (Q2) Cải thiện đó có còn khi kiểm tra ngoài phân phối (JSRT) và khi
che nốt không?
\end{keybox}

Ngưỡng 0,60 nằm cùng khoảng với sensitivity của người đọc trên các tập làm giàu ca bệnh nêu trên, nhưng không
mốc nào đo ở cùng điểm vận hành specificity 0,90. Vì vậy báo cáo \emph{không} so sánh ``ngang bác sĩ''; ngưỡng
này cũng thấp hơn rõ so với CAD chuyên biệt~\cite{vanleeuwen2024commercial,lee2020dl}.

\subsection{Đóng góp}

\begin{enumerate}
  \item \textbf{Benchmark nốt CXR chuẩn hoá cho VLM}, với prompt và parser công khai, chọn ngưỡng trên
        validation và ba tập test (trong phân phối, ngoài phân phối, có phân cấp độ mờ) --- bản tái lập có kiểm
        soát và quy mô lớn hơn của~\cite{khovanova2026comparative,nishio2026zeroshot}.
  \item \textbf{Nguyên mẫu trợ lý dựa trên VLM cỡ nhỏ tinh chỉnh LoRA cho phát hiện và định vị nốt/khối phổi},
        so trực tiếp với CNN classifier và detector trên cùng split; theo hiểu biết của nhóm, đây là bằng chứng
        đầu tiên cho tác vụ này trong phạm vi tài liệu đã khảo sát.
  \item \textbf{Phân tích chế độ báo bình thường và tác dụng của cân bằng lớp} khi tinh chỉnh VLM cho tổn
        thương hiếm và nhỏ --- phần gắn trực tiếp với nội dung học có mất cân bằng lớp của môn học.
  \item \textbf{Kiểm toán ``đúng vì đúng lý do''}: kết hợp IoU với phép che nốt so với che vùng ngẫu nhiên,
        phân tầng theo kích thước và độ mờ.
\end{enumerate}

\begin{warnbox}
\textbf{Phạm vi.} ``Trợ lý'' trong đồ án là \emph{nguyên mẫu nghiên cứu}. Kết quả không dùng cho quyết định
lâm sàng và không được trình bày như một công cụ sẵn sàng triển khai (xem Mục~\ref{sec:conclusion}).
\end{warnbox}

\subsection{Cấu trúc báo cáo}

Mục~\ref{sec:background} trình bày cơ sở lý thuyết và tổng quan tài liệu. Mục~\ref{sec:method} mô tả kiến trúc
trợ lý, các cấu hình và quy trình huấn luyện. Mục~\ref{sec:data} mô tả dữ liệu. Mục~\ref{sec:results} trình
bày thiết kế thí nghiệm, cách đánh giá và kết quả. Mục~\ref{sec:conclusion} kết luận, nêu hạn chế và hướng
phát triển.
